\section{Personalización de la interfaz gráfica}

La interfaz gráfica de Voctomix 2.0 se ha adaptado a los requisitos de una producción
de tipo evento técnico. Los cambios se concentran en tres áreas principales:

\subsection{Barra de herramientas}

Se han añadido los botones de Stream Blanker (LIVE, PAUSE, NOSTREAM), el selector de
overlay activo con vista previa en miniatura y el indicador de estado del notario.
El \textit{layout} de la barra se define en el fichero XML
\texttt{voctogui/ui/voctogui.ui}, que describe la jerarquía de widgets GTK3 y permite
reorganizar los elementos sin modificar el código Python.

\subsection{Panel de previsualizaciones}

El panel lateral muestra cuatro miniaturas JPEG de las fuentes activas (cam1--cam4),
actualizadas a 25 fps desde los puertos de previsualización 14000--14003. Cada
miniatura incluye un indicador de nivel de audio (\textit{VU-meter}) y un slider de
volumen individual para la monitorización de sala, independiente del volumen emitido
al stream.

\subsection{Reloj de estudio}

El módulo \texttt{voctogui/lib/studioclock.py} implementa un reloj de estudio que
muestra la hora actual y la duración de la sesión en curso, sincronizado con el
servidor mediante el reloj de red GStreamer. El reloj es visible en todo momento
en la esquina superior de la interfaz, permitiendo al realizador gestionar los tiempos
de la producción sin recurrir a herramientas externas.

% --- Figura sugerida ---
% \begin{figure}[H]
%   \centering
%   \includegraphics[width=0.9\textwidth]{figuras/gui_anotada.png}
%   \caption{Captura anotada de la interfaz gráfica de Voctomix 2.0.}
%   \label{fig:gui_anotada}
% \end{figure}
